% STATUS: written 2026-11-06 (English gate clean; sources read in full as recorded below).
% Sources read in full: Gerstenhaber 1964 (Ann. of Math. 79, 59--103), Gerstenhaber--Schack 1983 (JPAA 30, 143--156).
\section{Gerstenhaber structure and deformations}\label{sec:cohomology:gerstenhaber}

The first group $H^1(A,A)$ is the group of infinitesimal automorphisms of $A$,
that is the derivations modulo the inner
ones~\cite[pp.~59--60]{Gerstenhaber1964}. Gerstenhaber asked what the higher
groups do, and his 1964 paper answers for the second
one~\cite[pp.~59--60]{Gerstenhaber1964}. The group $H^2(A,A)$ is the group of
infinitesimal deformations
\cite[pp.~59--60]{Gerstenhaber1964}. It behaves toward deformations of $A$ as
$H^1$ behaves toward
automorphisms~\cite[pp.~59--60]{Gerstenhaber1964}. From this he reads a
rigidity criterion.

Gerstenhaber states the conclusion in his
introduction~\cite[p.~60]{Gerstenhaber1964}. The vanishing of $H^2(A,A)$ is a
sufficient condition for the rigidity of $A$~\cite[p.~60]{Gerstenhaber1964}.
Whether it is also necessary is open~\cite[p.~60]{Gerstenhaber1964}. It follows
at once that separable semi-simple algebras are
rigid~\cite[pp.~59--60]{Gerstenhaber1964}.

The deformation theory comes with
obstructions~\cite[\S 6]{Gerstenhaber1964}. An infinitesimal deformation need
not integrate to a true deformation~\cite[\S 6]{Gerstenhaber1964}. Gerstenhaber
defines a cohomology operation, the square, mapping $H^1(A,A)$ into $H^2(A,A)$
\cite[\S 6]{Gerstenhaber1964}. The obstruction to integrating an infinitesimal
deformation is a $3$-cocycle~\cite[pp.~60--61, \S 6]{Gerstenhaber1964}. He
develops this in his section on the additivity and integrability of the
square~\cite[pp.~60--61, \S 6]{Gerstenhaber1964}.

For fields the commutative theory gives a clean answer (Theorem~\ref{thm:cohomology:gerstenhaber3}).

\begin{theorem}[{\cite[Theorem 3, p.~78]{Gerstenhaber1964}}]\label{thm:cohomology:gerstenhaber3}
Let $K$ be an extension of a field $k$. Then $K$ is separable over $k$ if and only
if the second commutative cohomology group of $K$, considered as an algebra over
$k$, with coefficients in itself vanishes.
\end{theorem}

The higher Hochschild groups carry more structure than the groups alone
suggest~\cite[p.~144]{GerstenhaberSchack1983}. Gerstenhaber and Schack later
identified a graded Lie product on $H^*(A,A)$,
possibly non-abelian~\cite[p.~144]{GerstenhaberSchack1983}. They found that it had
gone unnoticed on the simplicial side, for a simple reason. Their main theorem
identifies the simplicial cohomology of a locally finite complex $C$ with the
Hochschild cohomology of the algebra $kC$. It also preserves the cup
product~ (Theorem~\ref{thm:cohomology:sc-hc})\cite[pp.~143--144]{GerstenhaberSchack1983}.

\begin{theorem}[{SC $=$ HC, \cite[pp.~143--144]{GerstenhaberSchack1983}}]\label{thm:cohomology:sc-hc}
Let $k$ be a commutative unital ring. There is a functor $C \mapsto kC$ from
locally finite simplicial complexes to associative unital $k$-algebras inducing an
isomorphism $H^*(C,k) \cong H^*(kC,kC)$ between the simplicial cohomology of $C$
and the Hochschild cohomology of $kC$, preserving the cup product.
\end{theorem}

When the algebra is $kC$ this graded Lie product is abelian, which explains why no
such product had been recorded on $H^*(C,k)$~\cite[p.~144]{GerstenhaberSchack1983}.

% Moved from chapters/20-applications/10-noncommutative-geometry.tex (sources read: Ginzburg--Schedler 2010, Sardanashvily 2007).
